\documentclass[11pt]{article}
\usepackage{amsmath}
\usepackage{amssymb}
\usepackage[margin=2cm]{geometry}
\usepackage{longtable}
\usepackage{booktabs}
\usepackage{pdflscape}
\usepackage{breqn}
\title{Certificate for [sqrt(2)]\_q (first 13 coefficients)}
\author{qreals certificate}
\date{\today}
\begin{document}
\maketitle

\noindent Input: the real number x = sqrt(2). The continued fraction is truncated to the depth that locks the first 13 coefficients (MGO Proposition 1.1).

\section*{(a) Continued fraction and even-length MGO form}
Regular continued fraction: [1, 2, 2, 2, 2, 2, 2, 2]. \par
Even-length MGO form: [1, 2, 2, 2, 2, 2, 2, 2]. \par
The even-length form evaluates to the convergent $ \frac{577}{408} $.\par

\section*{(b) MGO formula folded step by step}
Odd positions carry $[a]_q$ with $q^a$ above; even positions carry
$[a]_{q^{-1}}$ with $q^{-a}$ above (MGO eqn.~1.1).

\noindent\textbf{innermost term a\_8 = 2:}
{\footnotesize%
\begin{dmath*}
\frac{q + 1}{q}
\end{dmath*}
}

\noindent\textbf{fold in a\_7 = 2:}
{\footnotesize%
\begin{dmath*}
\frac{q^{3} + q^{2} + 2 q + 1}{q + 1}
\end{dmath*}
}

\noindent\textbf{fold in a\_6 = 2:}
{\footnotesize%
\begin{dmath*}
\frac{q^{5} + 2 q^{4} + 3 q^{3} + 3 q^{2} + 2 q + 1}{q^{5} + q^{4} + 2 q^{3} + q^{2}}
\end{dmath*}
}

\noindent\textbf{fold in a\_5 = 2:}
{\footnotesize%
\begin{dmath*}
\frac{q^{7} + 2 q^{6} + 5 q^{5} + 6 q^{4} + 6 q^{3} + 5 q^{2} + 3 q + 1}{q^{5} + 2 q^{4} + 3 q^{3} + 3 q^{2} + 2 q + 1}
\end{dmath*}
}

\noindent\textbf{fold in a\_4 = 2:}
{\footnotesize%
\begin{dmath*}
\frac{q^{9} + 3 q^{8} + 7 q^{7} + 11 q^{6} + 13 q^{5} + 13 q^{4} + 11 q^{3} + 7 q^{2} + 3 q + 1}{q^{9} + 2 q^{8} + 5 q^{7} + 6 q^{6} + 6 q^{5} + 5 q^{4} + 3 q^{3} + q^{2}}
\end{dmath*}
}

\noindent\textbf{fold in a\_3 = 2:}
{\footnotesize%
\begin{dmath*}
\frac{q^{11} + 3 q^{10} + 9 q^{9} + 16 q^{8} + 24 q^{7} + 29 q^{6} + 29 q^{5} + 25 q^{4} + 18 q^{3} + 10 q^{2} + 4 q + 1}{q^{9} + 3 q^{8} + 7 q^{7} + 11 q^{6} + 13 q^{5} + 13 q^{4} + 11 q^{3} + 7 q^{2} + 3 q + 1}
\end{dmath*}
}

\noindent\textbf{fold in a\_2 = 2:}
\noindent The value is the ratio $N(q)/D(q)$ with:\par
{\footnotesize%
\begin{dmath*}
N(q) = q^{13} + 4 q^{12} + 12 q^{11} + 25 q^{10} + 41 q^{9} + 56 q^{8} + 65 q^{7} + 65 q^{6} + 56 q^{5} + 41 q^{4} + 25 q^{3} + 12 q^{2} + 4 q + 1
\end{dmath*}
}

{\footnotesize%
\begin{dmath*}
D(q) = q^{13} + 3 q^{12} + 9 q^{11} + 16 q^{10} + 24 q^{9} + 29 q^{8} + 29 q^{7} + 25 q^{6} + 18 q^{5} + 10 q^{4} + 4 q^{3} + q^{2}
\end{dmath*}
}

\noindent\textbf{fold in a\_1 = 1:}
\noindent The value is the ratio $N(q)/D(q)$ with:\par
{\footnotesize%
\begin{dmath*}
N(q) = q^{14} + 4 q^{13} + 13 q^{12} + 28 q^{11} + 49 q^{10} + 70 q^{9} + 85 q^{8} + 90 q^{7} + 83 q^{6} + 66 q^{5} + 45 q^{4} + 26 q^{3} + 12 q^{2} + 4 q + 1
\end{dmath*}
}

{\footnotesize%
\begin{dmath*}
D(q) = q^{13} + 4 q^{12} + 12 q^{11} + 25 q^{10} + 41 q^{9} + 56 q^{8} + 65 q^{7} + 65 q^{6} + 56 q^{5} + 41 q^{4} + 25 q^{3} + 12 q^{2} + 4 q + 1
\end{dmath*}
}

\noindent\textbf{Result.} [577/408]\_q (the convergent, whose Taylor expansion gives [x]\_q):
\noindent The value is the ratio $N(q)/D(q)$ with:\par
{\footnotesize%
\begin{dmath*}
N(q) = q^{14} + 4 q^{13} + 13 q^{12} + 28 q^{11} + 49 q^{10} + 70 q^{9} + 85 q^{8} + 90 q^{7} + 83 q^{6} + 66 q^{5} + 45 q^{4} + 26 q^{3} + 12 q^{2} + 4 q + 1
\end{dmath*}
}

{\footnotesize%
\begin{dmath*}
D(q) = q^{13} + 4 q^{12} + 12 q^{11} + 25 q^{10} + 41 q^{9} + 56 q^{8} + 65 q^{7} + 65 q^{6} + 56 q^{5} + 41 q^{4} + 25 q^{3} + 12 q^{2} + 4 q + 1
\end{dmath*}
}

\noindent Taylor coefficients:
{\small%
\begin{dmath*}
1 + q^{3} - 2 q^{5} + q^{6} + 4 q^{7} - 5 q^{8} - 7 q^{9} + 18 q^{10} + 7 q^{11} - 55 q^{12} + O(q^{13})
\end{dmath*}
}

\section*{(c) Independent cross-checks}
\begin{itemize}
\item \textbf{[pass]} the first 13 coefficients are unchanged when 17 are asked for
\item \textbf{[pass]} [x+1]\_q computed from its own continued fraction equals q*[x]\_q + 1
\item \textbf{[na]} the series [sqrt(2)]\_q diverges at q=1; q=1 is for rationals
\item \textbf{[na]} sqrt(2) is irrational, so the continued fraction does not terminate
\end{itemize}
Summary: verified: truncation stable to 13, shift law [x+1]=q[x]+1; n/a: q=1 specialisation; exact rational function.

\section*{References}
S. Morier-Genoud and V. Ovsienko, q-deformed rationals and q-continued
fractions, Forum Math. Sigma \textbf{8} (2020), e13. Per-function theorem and
check mapping: \texttt{docs/CORRECTNESS.md}.

\medskip
\noindent This is a human-auditable derivation, not a formal machine proof;
every line above is checkable by hand.
\end{document}
